% ============================ 第三节 风险因素 ================================
% 分类对照准则第 57 号第二十三条：与发行人相关的风险 / 与行业相关的风险 /
% 其他风险；按重要性排序、结合公司实际定量披露，不含风险对策、竞争优势
% 及类似表述（第二十七条边界），一项风险因素不描述多个风险。
\gssection{风险因素}

投资者在评价发行人本次发行的股票时，除本招股说明书提供的其他资料外，应特别
认真地考虑下述各项风险因素。下述风险按照重要性原则或可能影响投资决策的程度
大小排序，该排序并不表示风险因素依次发生。

\subsection{与发行人相关的风险}

\subsubsection{技术升级与产品迭代风险}

人工智能大模型技术仍处于快速演进阶段，模型架构、训练方法、推理优化与产品形
态迭代频繁，研发投入强度高、试错成本大。如果公司对技术路线的判断出现偏差，
或者新一代模型的训练效果、研发进度不及预期，公司模型能力与产品体验可能落后
于同行业竞争者，进而对公司经营业绩产生不利影响。

\subsubsection{核心技术人员流失及技术失密风险}

截至2025年12月31日，公司拥有研发与技术人员726人，其中核心技术人员4名。人工
智能领域高层次人才供给稀缺、行业竞争白热化，人才薪酬水平持续上升。若公司核
心技术人员流失，或者出现模型权重、训练数据与核心算法泄密、被侵权等情形，将
对公司的技术研发和生产经营造成不利影响。

\subsubsection{知识产权及开源合规风险}

截至2025年12月31日，公司拥有境内授权专利236项、软件著作权120项。人工智能领
域的知识产权保护规则仍在发展完善过程中，模型训练数据、模型输出内容的权属界
定存在不确定性；同时，公司在研发过程中使用了部分开源软件及开源模型组件，若
未能持续满足相关开源协议要求，或者第三方对公司训练数据、模型成果主张权利，
可能引发纠纷并对公司经营产生不利影响。

\subsubsection{客户相对集中风险}

报告期各期，公司对前五大客户的销售收入分别为35,090.50万元、45,547.75万元和
58,232.93万元，占营业收入的比例分别为42.56\%、43.28\%和43.92\%，客户集中度
相对较高。虽然公司与主要客户建立了稳定的合作关系，但若主要客户因预算调整、
采购策略变化等原因减少对公司产品及服务的采购，公司存在经营业绩下滑的风险。

\subsubsection{应收账款回收风险}

报告期各期末，公司应收账款账面余额分别为18,236.45万元、23,845.67万元和
30,124.56万元，占当期营业收入的比例分别为22.12\%、22.66\%和22.72\%，随经营
规模扩大而增长。公司解决方案类业务客户以大型企业及政务客户为主、回款周期较
长，若主要客户信用状况发生不利变化，公司存在应收账款发生坏账或回收期延长的
风险。

\subsubsection{合同履约成本及存货余额较大的风险}

报告期各期末，公司存货账面余额分别为15,384.26万元、19,876.35万元和24,538.47
万元，主要为尚未终验收的私有化部署及解决方案项目的合同履约成本和外购软硬件。
若客户项目变更、取消或者验收进度延后，相关成本存在减值及占用营运资金的风险。

\subsubsection{毛利率波动风险}

报告期各期，公司综合毛利率分别为42.35\%、43.18\%和44.02\%。公司毛利率受业务
结构、算力成本、项目实施成本及市场竞争状况等多重因素影响，若开放平台服务价
格因竞争下调，或者算力等成本大幅上升，公司毛利率存在波动乃至下滑的风险。

\subsubsection{税收优惠政策变化风险}

公司为高新技术企业，报告期内按15\%的优惠税率缴纳企业所得税，并享受研发费用
加计扣除等税收优惠政策。若相关税收优惠政策发生变化，或者公司未来不能持续取
得高新技术企业资格，将对公司净利润水平产生不利影响。

\subsubsection{实际控制人控制不当的风险}

本次发行前，公司实际控制人林××合计控制公司63.00\%的表决权；本次发行完成后，
其控制的表决权比例仍达47.25\%，仍处于控制地位。若实际控制人利用其控制地位，
对公司经营决策、人事安排、利润分配等施加不当影响，可能损害公司及中小股东的
利益。

\subsubsection{经营规模扩大带来的管理风险}

本次发行上市后，公司资产规模、业务规模与人员规模将进一步扩大，对公司在战略
规划、组织设置、运营管理、内部控制等方面提出更高要求。若公司管理水平不能适
应规模扩张的需要，将影响公司的经营效率与发展质量。

\subsubsection{募集资金投资项目风险}

募集资金投资项目系基于当前技术发展趋势、市场环境与公司发展规划作出的安排。
项目实施过程中，若技术路线或市场环境发生重大不利变化、项目组织实施不力或者
进度不及预期，募投项目存在无法实现预期效益的风险。同时，募投项目实施后公司
固定资产及无形资产规模将明显增加，若收入增长不能同步覆盖新增折旧摊销，将对
公司经营业绩产生不利影响。

\subsection{与行业相关的风险}

\subsubsection{算力供应及核心芯片受限风险}

大模型的训练与推理高度依赖高性能智能计算芯片及配套算力基础设施，相关芯片的
供应受国际贸易环境及出口管制政策影响存在不确定性。若高性能芯片供应持续紧张、
采购或租赁成本大幅上升，或者国产替代方案的性能与生态成熟度不及预期，将对公
司模型迭代进度、服务交付能力及盈利能力产生不利影响。

\subsubsection{数据合规与内容安全风险}

生成式人工智能服务受《网络安全法》《数据安全法》《个人信息保护法》《生成式
人工智能服务管理暂行办法》等法律法规及监管规则约束，在训练数据来源合规、个
人信息保护、算法备案与生成内容安全等方面须持续满足监管要求。若公司未能持续
符合相关要求，或者生成内容出现违法违规情形，公司可能面临监管处罚、服务下架
及声誉受损的风险。

\subsubsection{市场竞争加剧风险}

人工智能大模型行业参与者众多，头部科技企业、专业大模型厂商及开源社区并存，
技术迭代与价格竞争激烈，开放平台服务价格呈下降趋势。若公司不能持续保持模型
能力、产品与服务的竞争力，市场竞争加剧可能导致公司市场份额下降、产品及服务
价格与毛利率下滑。

\subsubsection{行业技术路线变革风险}

人工智能基础理论与工程范式仍在快速发展，未来若出现更具效率或能力优势的新型
模型架构、训练范式，现有技术积累的相对价值可能下降，行业竞争格局可能重塑，
公司存在因技术路线变革而竞争地位下降的风险。

\subsubsection{下游需求及商业化进程不及预期风险}

大模型在各行业的应用仍处于渗透早期，下游客户的预算投入与应用深度受宏观经济
环境、行业景气度及应用效果影响。若行业整体商业化进程放缓或客户支出收缩，公
司业务增长存在放缓的风险。

\subsection{其他风险}

\subsubsection{发行失败风险}

本次发行结果将受到证券市场整体情况、投资者对公司价值判断等多种内外部因素的
影响。若出现认购不足等情形，本次发行存在发行失败的风险。

\subsubsection{不可抗力风险}

若发生自然灾害、公共卫生事件等不可抗力事件，可能影响公司及上下游的正常经营，
进而对公司经营业绩产生不利影响。
